\chapter{Boolean Analysis — Circuit Degree}
%%% generated by the blueprint pipeline from TCSlib.BooleanAnalysis.RazborovSmolensky.CircuitDegree
%%%%%%%%%%%%%%%%%%%%%%%%%%%%%%%%%%%%%%%%%%%%%%%%%%%%%%%%%%%%%%%%%%%%%%%%%%%%%%%
\section{Overview}

This file carries out the Razborov--Smolensky polynomial approximation of an
$\mathrm{AC}^0[p]$ circuit: layer by layer, every gate of a feed-forward circuit is replaced by a
random low-degree polynomial over $\bbz/p$, and the layerwise error probabilities are combined by a
union bound.  The end products are existence theorems producing a seed-indexed distribution (or an
explicit list) of polynomials of total degree at most $((p-1)\ell)^{\mathrm{depth}}$ that agree with
the circuit at any fixed input except on a small fraction of seeds.

%%%%%%%%%%%%%%%%%%%%%%%%%%%%%%%%%%%%%%%%%%%%%%%%%%%%%%%%%%%%%%%%%%%%%%%%%%%%%%%
\section{Declarations}

\begin{definition}[Boolean input cast to the prime field]
\lean{ACP.boolInput}
\label{ACP.boolInput}
\leanok
For a Boolean input $x : \mathrm{Fin}\,n \to \mathrm{Fin}\,2$, the map
$\mathrm{boolInput}_p(x)$ sends each coordinate $i$ to the image of the natural number $x_i$ in
$\bbz/p$.
\end{definition}

\begin{definition}[Field value of a bit]
\lean{ACP.boolVal}
\label{ACP.boolVal}
\leanok
The value in $\bbz/p$ represented by a bit $b : \mathrm{Fin}\,2$, namely the image of the natural
number $b$ under the canonical map $\N \to \bbz/p$.
\end{definition}

\begin{lemma}[Bit values lie in $\{0,1\}$]
\lean{ACP.boolVal_mem}
\label{ACP.boolVal_mem}
\leanok
\uses{ACP.boolVal}
\difficulty{1}
Let $p$ be a prime and let $b : \mathrm{Fin}\,2$ be a bit. Then the element of $\bbz/p$
represented by $b$ — the image of the natural number $b$ under the canonical map $\N \to
\bbz/p$ — belongs to the subset $\{0,1\}$ of $\bbz/p$.
\end{lemma}

\begin{lemma}[Booleanization recovers a bit from its field value]
\lean{ACP.bitify_boolVal}
\label{ACP.bitify_boolVal}
\leanok
\uses{ACP.bitify, ACP.boolVal}
\difficulty{1}
Let $p$ be a prime and let $b \in \mathrm{Fin}\,2$ be a bit. Booleanizing the field
value of $b$ recovers $b$: taking the element of $\bbz/p$ represented by $b$ and then
mapping it back to $\mathrm{Fin}\,2$ (sending $1$ to $1$ and everything else to $0$)
yields $b$ again.
\end{lemma}

\begin{lemma}[One-step unfolding of node evaluation]
\lean{ACP.evalNode_succ_eq}
\label{ACP.evalNode_succ_eq}
\leanok
\uses{ACP.FeedForward, ACP.FeedForward.evalNode}
\difficulty{1}
Let $F$ be a feedforward circuit over an alphabet $\alpha$, let $d$ be a layer index
with $0 \le d < \mathrm{depth}(F)$, let $u$ be a node of layer $d+1$, and let $x$ be an
assignment of alphabet values to the input nodes of $F$. Writing $g$ for the gate
attached to $u$, the value of $u$ under $x$ equals the gate operation of $g$ applied to
the tuple whose entry at each input slot $i$ is the value under $x$ of the layer-$d$
node to which slot $i$ is wired.
\end{lemma}

\begin{definition}[Degree target after $d$ layers]
\lean{ACP.circuitDegreeBound}
\label{ACP.circuitDegreeBound}
\leanok
The degree budget for the approximating polynomials after $d$ layers, defined as
\[
\mathrm{circuitDegreeBound}(p,\ell,d) = \bigl((p-1)\ell\bigr)^{d}.
\]
\end{definition}

\begin{definition}[Gate count in the first $d$ layers]
\lean{ACP.gateCountBefore}
\label{ACP.gateCountBefore}
\leanok
\uses{ACP.FeedForward}
For a circuit $F$ with finite node sets, $\mathrm{gateCountBefore}\,F\,d$ counts the non-input gates
in layers $1,\dots,d$: it is $0$ for $d = 0$ and, recursively, the count for $d$ plus the number of
nodes at layer $d+1$.
\end{definition}

\begin{lemma}[Vanishing of the gate count through layer zero]
\lean{ACP.gateCountBefore_zero}
\label{ACP.gateCountBefore_zero}
\leanok
\uses{ACP.FeedForward, ACP.gateCountBefore}
\difficulty{1}
Let $F$ be a layered feedforward circuit over the Boolean alphabet $\{0,1\}$, whose
input layer carries $n$ nodes and each of whose layers has a finite node set. Then the
number of non-input gates lying in the first $0$ layers is zero, since the input layer
contributes no gates.
\end{lemma}

\begin{lemma}[Gate-count recursion across layers]
\lean{ACP.gateCountBefore_succ}
\label{ACP.gateCountBefore_succ}
\leanok
\uses{ACP.FeedForward, ACP.gateCountBefore}
\difficulty{1}
Let $F$ be a layered feedforward circuit over the Boolean alphabet $\bbf_2$ with input
type $\mathrm{Fin}\,n$, in which every layer has finitely many nodes, and for each
natural number $k \le \mathrm{depth}(F)$ let $g_F(k)$ denote the number of non-input
gates lying in layers $1,\dots,k$ (so $g_F(0) = 0$). Then for every $d$ with $d + 1 \le
\mathrm{depth}(F)$, \[ g_F(d+1) = g_F(d) + \abs{N_{d+1}}, \] where $\abs{N_{d+1}}$ is
the number of nodes in layer $d+1$.
\end{lemma}

\begin{lemma}[Cardinality of a product filtered on the left factor]
\lean{ACP.prod_left_filter_card}
\label{ACP.prod_left_filter_card}
\leanok
\difficulty{4}
Let $\alpha$ and $\beta$ be finite types, and let $P$ be a decidable predicate on
$\alpha$. Then the number of pairs $(a,b) \in \alpha \times \beta$ whose first component
satisfies $P$ equals the number of elements $a \in \alpha$ satisfying $P$, multiplied by
the cardinality $\abs{\beta}$.
\end{lemma}

\begin{lemma}[Fiberwise product counting bound]
\lean{ACP.prod_filter_fiber_mul_le}
\label{ACP.prod_filter_fiber_mul_le}
\leanok
\difficulty{5}
Let $\alpha$ and $\beta$ be finite types, let $P$ be a predicate on $\alpha$ and $Q$ a
predicate on $\alpha \times \beta$, and let $C, B$ be natural numbers. Suppose that for
every $a \in \alpha$ satisfying $P(a)$, the fiber $\{\,b \in \beta : Q(a,b)\,\}$
satisfies $\abs{\{\,b : Q(a,b)\,\}} \cdot C \le B$. Then
\[
\abs{\{\,(a,b) \in \alpha \times \beta : P(a) \wedge Q(a,b)\,\}} \cdot C
\;\le\;
\abs{\{\,a \in \alpha : P(a)\,\}} \cdot B .
\]
\end{lemma}

\begin{definition}[Splitting off one coordinate of a dependent product]
\lean{ACP.piEquivAt}
\label{ACP.piEquivAt}
\leanok
For an index $i$ with decidable equality on the index type, the equivalence
\[
\bigl((j : \iota) \to \beta_j\bigr) \;\simeq\; \beta_i \times \bigl((j : \{j : j \ne i\}) \to \beta_j\bigr)
\]
sending a dependent function to its value at $i$ together with its restriction to the remaining
coordinates.
\end{definition}

\begin{lemma}[Counting functions bad at a fixed coordinate]
\lean{ACP.pi_coordinate_bad_mul_le}
\label{ACP.pi_coordinate_bad_mul_le}
\leanok
\uses{ACP.piEquivAt}
\difficulty{5}
Let $\iota$ be a finite index type with decidable equality, let $(\beta_j)_{j \in
\iota}$ be a family of finite types, and fix a coordinate $i \in \iota$. Let
$\mathrm{Bad}$ be a decidable predicate on $\beta_i$, and let $C$ be a natural number
such that
\[
\abs{\{\, b \in \beta_i : \mathrm{Bad}(b) \,\}} \cdot C \;\le\; \abs{\beta_i}.
\]
Then the number of dependent functions $f \in \prod_{j \in \iota} \beta_j$ whose value
at $i$ is bad satisfies
\[
\abs{\{\, f \in \textstyle\prod_{j \in \iota} \beta_j : \mathrm{Bad}(f(i)) \,\}} \cdot C
\;\le\; \abs{\textstyle\prod_{j \in \iota} \beta_j}.
\]
\end{lemma}

\begin{lemma}[Counting product tuples bad in some coordinate]
\lean{ACP.pi_exists_bad_card_mul_le}
\label{ACP.pi_exists_bad_card_mul_le}
\leanok
\uses{ACP.pi_coordinate_bad_mul_le}
\difficulty{4}
Let $\iota$ be a finite index set, and for each $i \in \iota$ let $\beta_i$ be a finite
set equipped with a predicate $\mathrm{Bad}_i$ singling out its "bad" elements. Fix a
natural number $C$, and suppose that for every coordinate $i$ the bad elements are
sufficiently rare, in the sense that
\[
\abs{\{\, b \in \beta_i : \mathrm{Bad}_i(b) \,\}}\cdot C \;\le\; \abs{\beta_i}.
\]
Then, among all dependent tuples $f$ in the product $\prod_{i \in \iota} \beta_i$, those
that are bad in at least one coordinate satisfy
\[
\abs{\{\, f \in \textstyle\prod_{i \in \iota}\beta_i : \exists\, i,\ \mathrm{Bad}_i(f_i)
\,\}}\cdot C \;\le\; \abs{\iota}\cdot\Bigl\lvert\textstyle\prod_{i \in
\iota}\beta_i\Bigr\rvert.
\]
\end{lemma}

\begin{lemma}[Length is preserved when filtering commutes with mapping]
\lean{ACP.list_filter_map_length}
\label{ACP.list_filter_map_length}
\leanok
\difficulty{2}
Let $l$ be a finite list with entries in a set $\alpha$, let $f\colon \alpha \to \beta$
be a function, and let $P$ be a decidable predicate on $\beta$. Then applying $f$ to
every entry of $l$ and retaining those whose image satisfies $P$ yields a list of the
same length as retaining those entries $a$ of $l$ for which $P(f(a))$ holds; that is,
the sublist $\bigl(\mathrm{map}\,f\,(l)\bigr)$ filtered by $P$ and the sublist of $l$
filtered by the pulled-back predicate $a \mapsto P(f(a))$ have equal length.
\end{lemma}

\begin{lemma}[Filtered list length equals filtered finset cardinality]
\lean{ACP.finset_toList_filter_length_eq_card}
\label{ACP.finset_toList_filter_length_eq_card}
\leanok
\difficulty{3}
Let $s$ be a finite set of elements of a type $\alpha$, and let $q$ be a decidable
predicate on $\alpha$. Writing the elements of $s$ out as a list (in any order) and
discarding those that fail $q$ produces a list whose length equals the number of
elements of $s$ that satisfy $q$; that is, the length of the filtered list agrees with
the cardinality of the subset $\{x \in s : q(x)\}$.
\end{lemma}

\begin{definition}[Gate approximator family]
\lean{ACP.GatePolyFamily}
\label{ACP.GatePolyFamily}
\leanok
\uses{ACP.bitify, ACP.boolInput}
A \texttt{GatePolyFamily} for a gate operation $\mathrm{op}$ packages a nonempty finite seed type
together with, for each tuple of incoming polynomials and each seed, an approximating polynomial in
$n$ variables over $\bbz/p$ subject to two guarantees: its total degree is at most
$(p-1)\ell \cdot \sup_i \deg(\text{polys}_i)$, and for every Boolean input $x$ at which all incoming
polynomials evaluate into $\{0,1\}$, the number of seeds on which the approximator disagrees with
the true gate value, multiplied by $2^{\ell}$, is at most the number of seeds.  Crucially the seed
type depends only on the gate, not on the incoming polynomials.
\end{definition}

\begin{lemma}[Existence of a gate approximator family]
\lean{ACP.exists_gate_poly_family}
\label{ACP.exists_gate_poly_family}
\leanok
\uses{ACP.ACp_GateOps, ACP.ACp_GateOps_cases, ACP.GatePolyFamily, ACP.approxAnd, ACP.approxAnd_pointwise_bad_count, ACP.approxAnd_totalDegree, ACP.approxSeed_card, ACP.bitify, ACP.boolInput, ACP.cast_bitify_eq, ACP.exactAnd_on_bits, ACP.exactMod, ACP.exactMod_on_bits, ACP.exactMod_totalDegree, ACP.modGateOp}
\difficulty{6}
Let $p$ be a prime, and fix natural numbers $n$ and $\ell$. For every gate operation
$\mathrm{op}$ belonging to the $\mathrm{AC}^0[p]$ gate set — the identity, the NOT gate,
the unbounded fan-in AND gate of any arity, and the unbounded $\mathrm{MOD}_p$ gate of
any arity — there exists a gate approximator family for $\mathrm{op}$ in $n$ variables
with parameter $\ell$; that is, a nonempty finite seed type together with, for each
tuple of incoming polynomials over $\bbz/p$ and each seed, an approximating polynomial
in $n$ variables whose total degree is at most $(p-1)\,\ell$ times the largest degree
among the incoming polynomials, and such that for every Boolean input at which all
incoming polynomials evaluate into $\{0,1\}$, the number of seeds on which the
approximator disagrees with the true gate value, multiplied by $2^{\ell}$, is at most
the total number of seeds.
\end{lemma}

\begin{definition}[Chosen gate approximator family]
\lean{ACP.gatePolyFamily}
\label{ACP.gatePolyFamily}
\leanok
\uses{ACP.ACp_GateOps, ACP.GatePolyFamily, ACP.exists_gate_poly_family}
A choice of \texttt{GatePolyFamily} for each $\mathrm{AC}^0[p]$ gate operation, obtained from the
preceding existence statement.
\end{definition}

\begin{definition}[Layer polynomial family]
\lean{ACP.LayerPolyFamily}
\label{ACP.LayerPolyFamily}
\leanok
\uses{ACP.FeedForward, ACP.FeedForward.evalNode, ACP.boolInput, ACP.circuitDegreeBound, ACP.gateCountBefore}
A \texttt{LayerPolyFamily} for a circuit $F$ at layer $d$ consists of a nonempty finite seed type and,
for each seed, a polynomial for every node of layer $d$, such that each polynomial has total degree at
most $\mathrm{circuitDegreeBound}(p,\ell,d)$ and, for every Boolean input $x$, the number of seeds
for which some layer-$d$ node is mispredicted, multiplied by $2^{\ell}$, is at most
$\mathrm{gateCountBefore}\,F\,d$ times the number of seeds.
\end{definition}

\begin{definition}[Input layer family]
\lean{ACP.inputLayerFamily}
\label{ACP.inputLayerFamily}
\leanok
\uses{ACP.FeedForward, ACP.FeedForward.evalNode, ACP.LayerPolyFamily, ACP.boolInput, ACP.circuitDegreeBound, ACP.gateCountBefore}
The \texttt{LayerPolyFamily} at layer $0$: a single seed, with each input node represented exactly by
the corresponding variable $X_i$, so that no seed is ever bad.
\end{definition}

\begin{definition}[Inductive layer step]
\lean{ACP.stepLayerFamily}
\label{ACP.stepLayerFamily}
\leanok
\uses{ACP.ACp_GateOps, ACP.FeedForward, ACP.FeedForward.evalNode, ACP.FeedForward.onlyUsesGates, ACP.GatePolyFamily, ACP.LayerPolyFamily, ACP.bitify, ACP.bitify_boolVal, ACP.boolInput, ACP.boolVal_mem, ACP.circuitDegreeBound, ACP.evalNode_succ_eq, ACP.gateCountBefore, ACP.gatePolyFamily, ACP.pi_exists_bad_card_mul_le, ACP.prod_filter_fiber_mul_le, ACP.prod_left_filter_card}
Given a circuit $F$ using only $\mathrm{AC}^0[p]$ gates and a \texttt{LayerPolyFamily} at layer $d$,
this produces a \texttt{LayerPolyFamily} at layer $d+1$.  Its seed type is the product of the previous
seed type with one gate seed per node of layer $d+1$, and each node polynomial is the gate
approximator applied to the layer-$d$ polynomials of its inputs.
\end{definition}

\begin{definition}[Layer family built by recursion on depth]
\lean{ACP.buildLayerFamily}
\label{ACP.buildLayerFamily}
\leanok
\uses{ACP.ACp_GateOps, ACP.FeedForward, ACP.FeedForward.onlyUsesGates, ACP.LayerPolyFamily, ACP.inputLayerFamily, ACP.stepLayerFamily}
For a circuit using only $\mathrm{AC}^0[p]$ gates, the \texttt{LayerPolyFamily} at any layer
$d \le F.\mathrm{depth}$, obtained by starting from the input layer family and iterating the layer
step $d$ times.
\end{definition}

\begin{theorem}[{Polynomial approximation of the outputs of an AC⁰[p] circuit}]
\lean{ACP.exists_poly_distribution_for_circuit_outputs}
\label{ACP.exists_poly_distribution_for_circuit_outputs}
\leanok
\uses{ACP.ACp_GateOps, ACP.FeedForward, ACP.FeedForward.eval, ACP.FeedForward.evalNode, ACP.FeedForward.onlyUsesGates, ACP.boolInput, ACP.buildLayerFamily, ACP.circuitDegreeBound, ACP.gateCountBefore}
\difficulty{4}
Let $p$ be a prime, let $F$ be a feedforward circuit over the alphabet $\mathrm{Fin}\,2$
with input type $\mathrm{Fin}\,n$ and a finite output type, all of whose node sets are
finite and which uses only $\mathrm{AC}^0[p]$ gates, and fix an error parameter $\ell
\in \bbn$. Write $d = \mathrm{depth}(F)$ for its depth and $N$ for the total number of
its non-input gates, that is, the sum over the layers $1,\dots,d$ of the number of nodes
in each layer. For a bit $b \in \mathrm{Fin}\,2$ write $\overline{b} \in \bbz/p$ for the
image of $b$ regarded as a natural number, and for a Boolean input $x : \mathrm{Fin}\,n
\to \mathrm{Fin}\,2$ write $\hat{x} \in (\bbz/p)^n$ for the point whose $i$-th
coordinate is $\overline{x_i}$. Then there exist a nonempty finite seed type $S$ and,
for each seed $s \in S$ and each output node $o$, a polynomial $P_{s,o}$ in $n$
variables over $\bbz/p$, such that every $P_{s,o}$ has total degree at most
$\bigl((p-1)\ell\bigr)^{d}$ and, for every Boolean input $x$,
\[
\abs{\bigl\{\, s \in S : \exists\, o,\ P_{s,o}(\hat{x}) \ne \overline{F(x)(o)}
\,\bigr\}} \cdot 2^{\ell} \;\le\; N \cdot \abs{S},
\]
where $F(x)(o) \in \mathrm{Fin}\,2$ is the value the circuit $F$ computes at output node
$o$ on input $x$.
\end{theorem}

\begin{theorem}[{Polynomial approximation of a single-output AC⁰[p] circuit}]
\lean{ACP.exists_poly_distribution_for_circuit_one}
\label{ACP.exists_poly_distribution_for_circuit_one}
\leanok
\uses{ACP.ACp_GateOps, ACP.FeedForward, ACP.FeedForward.eval, ACP.FeedForward.evalNode, ACP.FeedForward.eval₁, ACP.FeedForward.onlyUsesGates, ACP.boolInput, ACP.buildLayerFamily, ACP.circuitDegreeBound, ACP.gateCountBefore}
\difficulty{4}
Let $p$ be prime and let $n, \ell \in \bbn$. Let $F$ be a feedforward circuit over the
alphabet $\mathrm{Fin}\,2$ with input type $\mathrm{Fin}\,n$, a single output node, and
a finite node set at every layer, each of whose gates is an $\mathrm{AC}^0[p]$ gate (the
identity, NOT, an unbounded fan-in AND, or an unbounded $\mathrm{MOD}_p$ gate of some
arity); write $d$ for its depth. For a Boolean input $x : \mathrm{Fin}\,n \to
\mathrm{Fin}\,2$, let $\hat{x}$ be its coordinatewise image in $\bbz/p$ and let $F(x)
\in \{0,1\} \subseteq \bbz/p$ be the resulting single output value cast into the field.
Then there exist a nonempty finite seed set $\mathrm{Seed}$ and a family of polynomials
$P_s \in (\bbz/p)[X_0,\dots,X_{n-1}]$ indexed by $s \in \mathrm{Seed}$, each of total
degree at most $\bigl((p-1)\ell\bigr)^{d}$, such that for every Boolean input $x$,
\[
\abs{\{\,s \in \mathrm{Seed} : P_s(\hat{x}) \ne F(x)\,\}} \cdot 2^{\ell}
\;\le\;
G \cdot \abs{\mathrm{Seed}},
\]
where $G$ is the total number of non-input gates of $F$, that is, the total number of
nodes across layers $1,\dots,d$.
\end{theorem}

\begin{theorem}[{Polynomial list approximating a single-output $\mathrm{AC}^0[p]$ circuit}]
\lean{ACP.exists_poly_list_for_circuit_one}
\label{ACP.exists_poly_list_for_circuit_one}
\leanok
\uses{ACP.ACp_GateOps, ACP.FeedForward, ACP.FeedForward.eval₁, ACP.FeedForward.onlyUsesGates, ACP.boolInput, ACP.circuitDegreeBound, ACP.exists_poly_distribution_for_circuit_one, ACP.finset_toList_filter_length_eq_card, ACP.gateCountBefore, ACP.list_filter_map_length}
\difficulty{4}
Let $p$ be a prime and let $F$ be a layered feedforward circuit over the Boolean
alphabet $\mathrm{Fin}\,2$, with inputs indexed by $\mathrm{Fin}\,n$, finitely many
nodes in every layer, and a single output node, and suppose every gate of $F$ belongs to
the $\mathrm{AC}^0[p]$ gate set (the identity, NOT, unbounded fan-in AND, and unbounded
$\mathrm{MOD}_p$ gates of every arity). Then for every $\ell \in \bbn$ there is a
nonempty list $P_1,\dots,P_m$ of polynomials in the variables $\{X_i : i \in
\mathrm{Fin}\,n\}$ over $\bbz/p$, each of total degree at most
$\bigl((p-1)\ell\bigr)^{\mathrm{depth}(F)}$, with the following property. For a Boolean
input $x : \mathrm{Fin}\,n \to \mathrm{Fin}\,2$, write $\hat{x} \in (\bbz/p)^n$ for the
point whose $i$-th coordinate is the image of $x_i$ in $\bbz/p$, and let $b(x) \in
\bbz/p$ be the image of the single Boolean value output by $F$ on $x$. Then for every
such $x$, the number of indices $j$ with $P_j(\hat{x}) \neq b(x)$, multiplied by
$2^{\ell}$, is at most $s \cdot m$, where $s$ is the total number of non-input gates of
$F$ and $m$ is the length of the list.
\end{theorem}
